% Comprehensive Paleomagnetic Results Tables
% Can be inserted into manuscript as needed

% ============================================
% TABLE 1: POLARITY CONFIRMATION RESULTS
% ============================================

\begin{table*}[htbp]
    \centering
    \caption{Paleomagnetic polarity confirmations for in-situ volcanic sites in the Bolaven Plateau. Ages from \textsuperscript{40}Ar/\textsuperscript{39}Ar dating by \citet{Sieh2019}. Expected polarities based on wes2020 geomagnetic polarity timescale model.}
    \label{tab:polarity_confirmations}
    \footnotesize
    \begin{tabular}{|l|c|c|l|l|c|c|c|c|c|c|}
        \hline
        \textbf{Site} & \textbf{Age (ka)} & \textbf{Age} & \textbf{Mag} & \textbf{Expected} & \textbf{Measured} & \textbf{Dec} & \textbf{Inc} & \textbf{K} & \textbf{$\alpha_{95}$} & \textbf{Match} \\
        & & \textbf{$\sigma$} & \textbf{Chron} & \textbf{Polarity} & \textbf{Polarity} & \textbf{(°)} & \textbf{(°)} & & \textbf{(°)} & \textbf{Status} \\
        \hline
        NWF-A & 518 & 6 & C1n & N & N & 10.3 & 21.2 & 79 & 4.3 & \checkmark \\
        NWF-B & 936 & 30? & C1r.1r & R & - & 81.7 & -60.8 & 5 & 16.9 & Clasts \\
        PKS-B & 906 & 12 & C1r.1r & R & R & 189.8 & -10.2 & 146 & 5.0 & \checkmark \\
        PKS-H & 770 & 9.3 & C1n/C1r.1r & N & R N N & 180.1 & -6.0 & 4 & 16.3 & Ejecta \\
        PKS-I & 553 & 20 & C1n & N & ? & 48.7 & -34.9 & 232 & 5.0 & Clast \\
        PKS-J & 601 & 44 & C1n & N & R & 175.6 & -32.5 & 77 & 5.5 & \textsuperscript{a} \\
        PKS-K & 407 & 18 & C1n & N & R & 203.3 & -33.5 & 18 & 12.6 & \textsuperscript{a} \\
        WBE-A & 1397 & 13 & C1r.3r & R & R & 167.0 & -33.8 & 79 & 5.2 & \checkmark \\
        \hline
    \end{tabular}
    
    \vspace{0.5em}
    \raggedright
    \footnotesize
    \textsuperscript{a}Assigned age likely incorrect; reversed polarity suggests Matuyama chron (>773 ka).\\
    Dec = declination, Inc = inclination, K = Fisher precision parameter, $\alpha_{95}$ = 95\% confidence angle.\\
    N = normal polarity, R = reversed polarity.
\end{table*}

% ============================================
% TABLE 2: CONGLOMERATE TEST RESULTS
% ============================================

\begin{table*}[htbp]
    \centering
    \caption{Watson (1956) conglomerate test results for suspected ejecta deposits in Bolaven diamicton. Random distribution (R < Ro\textsubscript{95}) indicates ballistic emplacement; systematic clustering (R > Ro\textsubscript{95}) indicates in-situ origin or post-depositional remagnetization.}
    \label{tab:conglomerate_tests}
    \footnotesize
    \begin{tabular}{|l|c|c|c|c|c|c|c|c|c|}
        \hline
        \textbf{Site} & \textbf{Site Type} & \textbf{N} & \textbf{R} & \textbf{Ro\textsubscript{95}} & \textbf{Ro\textsubscript{99}} & \textbf{K} & \textbf{$\alpha_{95}$} & \textbf{Test} & \textbf{Interpretation} \\
        & & & & & & & \textbf{(°)} & \textbf{Result} & \\
        \hline
        NWF-C & Diamicton & 8 & 3.37 & 4.48 & 5.26 & 1.5 & 74.5 & PASS & Ejecta \\
        PKQ-A & Double diamicton & 4 & 2.80 & 3.18 & 3.65 & 2.5 & 74.5 & PASS\textsuperscript{a} & Ejecta \\
        PKS-F & Coffee farm & 10 & 4.72 & 5.03 & 5.94 & 1.7 & 56.0 & PASS & Ejecta \\
        PKS-G & Coffee farm & 8 & 6.68 & 4.48 & 5.26 & 5.3 & 26.5 & FAIL & In-situ flow \\
        \hline
    \end{tabular}
    
    \vspace{0.5em}
    \raggedright
    \footnotesize
    \textsuperscript{a}Close to failure threshold but passes at 95\% confidence level.\\
    N = number of samples, R = resultant vector length, Ro = critical value for randomness,\\
    K = Fisher precision parameter, $\alpha_{95}$ = 95\% confidence angle.\\
    Pass: R < Ro\textsubscript{95}; Fail: R > Ro\textsubscript{95}.
\end{table*}

% ============================================
% TABLE 3: DETAILED SITE STATISTICS
% ============================================

\begin{table*}[htbp]
    \centering
    \caption{Complete Fisher statistics for all analyzed sites in the Bolaven Plateau paleomagnetic study. Sites classified as in-situ flows, ejecta deposits, or ambiguous based on polarity and conglomerate test results.}
    \label{tab:site_statistics}
    \footnotesize
    \begin{tabular}{|l|c|c|c|c|c|c|c|c|l|}
        \hline
        \textbf{Site} & \textbf{N\textsubscript{samp}} & \textbf{Dec} & \textbf{Inc} & \textbf{$\alpha_{95}$} & \textbf{K} & \textbf{R} & \textbf{n\textsubscript{lines}} & \textbf{n\textsubscript{planes}} & \textbf{Classification} \\
        & & \textbf{(°)} & \textbf{(°)} & \textbf{(°)} & & & & & \\
        \hline
        \multicolumn{10}{|c|}{\textit{In-Situ Basalt Flows}} \\
        \hline
        NWF-A & 5 & 10.3 & 21.2 & 4.3 & 79 & 14.8225 & 15 & 0 & In-situ (confirmed) \\
        NWF-B & 3 & 81.7 & -60.8 & 16.9 & 5 & 15.9602 & 20 & 0 & Transported? \\
        PKS-B & 1 & 189.8 & -10.2 & 5.0 & 146 & 6.9589 & 7 & 0 & In-situ (confirmed) \\
        PKS-I & 1 & 48.7 & -34.9 & 5.0 & 232 & 4.9828 & 5 & 0 & Clast (likely) \\
        PKS-J & 3 & 175.6 & -32.5 & 5.5 & 77 & 9.8830 & 10 & 0 & In-situ (age wrong) \\
        PKS-K & 2 & 203.3 & -33.5 & 12.6 & 18 & 8.5493 & 9 & 0 & In-situ (age wrong) \\
        WBE-A & 2 & 167.0 & -33.8 & 5.2 & 79 & 10.8735 & 11 & 0 & In-situ (confirmed) \\
        \hline
        \multicolumn{10}{|c|}{\textit{Ejecta/Diamicton Sites}} \\
        \hline
        NWF-C & 8 & 1.8 & 19.5 & 74.5 & 1.5 & 1.9925 & 2\textsuperscript{a} & 0 & Ejecta (confirmed) \\
        PKQ-A & 4 & 156.7 & 67.8 & 74.5 & 2.5 & - & - & 0 & Ejecta (confirmed) \\
        PKS-F & 10 & 28.6 & -12.9 & 56.0 & 1.7 & - & - & 0 & Ejecta (confirmed) \\
        PKS-G & 8 & 157.0 & 4.3 & 26.5 & 5.3 & - & - & 0 & In-situ flow \\
        PKS-H & 3 & 180.1 & -6.0 & 16.3 & 4 & 18.6149 & 24 & 0 & Ejecta (reinterpreted) \\
        \hline
    \end{tabular}
    
    \vspace{0.5em}
    \raggedright
    \footnotesize
    \textsuperscript{a}Average per sample.\\
    N\textsubscript{samp} = number of block/core samples, Dec = declination, Inc = inclination,\\
    $\alpha_{95}$ = 95\% confidence angle, K = Fisher precision parameter, R = resultant vector length,\\
    n\textsubscript{lines} = number of line fits, n\textsubscript{planes} = number of plane fits.
\end{table*}

% ============================================
% TABLE 4: INDIVIDUAL BOULDER/CORE DETAILS (NWF-C)
% ============================================

\begin{table*}[htbp]
    \centering
    \caption{Individual sample directions for site NWF-C demonstrating random distribution consistent with ejecta interpretation. Wide range of declinations and inclinations with diverse K values indicate preservation of pre-impact magnetizations.}
    \label{tab:nwfc_details}
    \footnotesize
    \begin{tabular}{|l|c|c|c|c|c|}
        \hline
        \textbf{Sample} & \textbf{Dec (°)} & \textbf{Inc (°)} & \textbf{K} & \textbf{Hemisphere} & \textbf{Notes} \\
        \hline
        NWFC-1 & 6.4 & 6.4 & 132 & Upper & Low angle, northern \\
        NWFC-2 & 345.0 & -43.7 & 298 & Lower & Steep negative inc \\
        NWFC-3 & 336.7 & 39.2 & 373 & Upper & NW direction \\
        NWFC-4 & 353.1 & -12.3 & 430 & Lower & Near north \\
        NWFC-5 & 130.2 & 22.4 & 285 & Upper & SE direction \\
        NWFC-6 & 61.6 & 53.8 & 85 & Upper & Steep positive \\
        NWFC-7 & 26.1 & 26.1 & 37 & Upper & NE direction \\
        NWFC-8 & 171.0 & 75.5 & 280 & Upper & Very steep \\
        \hline
    \end{tabular}
    
    \vspace{0.5em}
    \raggedright
    \footnotesize
    Range of declinations: 336.7° to 171.0° (full 360° coverage)\\
    Range of inclinations: -43.7° to 75.5° (119° spread)\\
    Both hemispheres represented (upper and lower)\\
    Diverse K values indicate independent magnetization histories
\end{table*}

% ============================================
% TABLE 5: PKS-H INDIVIDUAL SAMPLE RESULTS
% ============================================

\begin{table*}[htbp]
    \centering
    \caption{Individual core results for site PKS-H showing mixed polarities and scattered directions inconsistent with in-situ flow interpretation. HA1 and HA2 from same stratigraphic level show opposing polarities, supporting ejecta reinterpretation.}
    \label{tab:pksh_details}
    \footnotesize
    \begin{tabular}{|l|c|c|c|c|c|c|c|l|}
        \hline
        \textbf{Core} & \textbf{Dec (°)} & \textbf{Inc (°)} & \textbf{K} & \textbf{R} & \textbf{$\alpha_{95}$ (°)} & \textbf{n\textsubscript{lines}} & \textbf{Polarity} & \textbf{Notes} \\
        \hline
        HA1 & 141.0 & -29.0 & 59 & 9.8463 & 6.4 & 10 & R & Potential self-reversal? \\
        HA2 & 193.9 & 11.2 & 88 & 8.9095 & 5.5 & 9 & N & Same horizon as HA1 \\
        HB1 & 225.2 & 13.2 & 28 & 4.8590 & 14.6 & 5 & N & Below suspected boundary \\
        \hline
        \multicolumn{9}{|l|}{\textit{Site Mean (not geologically meaningful for ejecta)}} \\
        \hline
        All & 180.1 & -6.0 & 4 & 18.6149 & 16.3 & 24 & Mixed & Random distribution \\
        \hline
    \end{tabular}
    
    \vspace{0.5em}
    \raggedright
    \footnotesize
    HA1 and HA2 from same stratigraphic position but opposite polarities\\
    Low site mean K (4) and high $\alpha_{95}$ (16.3°) indicate poor clustering\\
    Scattered directions inconsistent with coherent in-situ flow\\
    \textbf{Interpretation: Ejecta deposit, not in-situ flow near BM boundary}
\end{table*}

% ============================================
% TABLE 6: AGE DISCREPANCIES
% ============================================

\begin{table*}[htbp]
    \centering
    \caption{Sites with polarity-age discrepancies requiring geochronological reassessment. Both sites show well-defined reversed polarities inconsistent with assigned Brunhes ages, suggesting actual ages in Matuyama chron (>773 ka).}
    \label{tab:age_problems}
    \footnotesize
    \begin{tabular}{|l|c|c|c|c|c|c|c|l|}
        \hline
        \textbf{Site} & \textbf{Assigned} & \textbf{Expected} & \textbf{Measured} & \textbf{Dec} & \textbf{Inc} & \textbf{K} & \textbf{$\alpha_{95}$} & \textbf{Suggested} \\
        & \textbf{Age (ka)} & \textbf{Polarity} & \textbf{Polarity} & \textbf{(°)} & \textbf{(°)} & & \textbf{(°)} & \textbf{True Age} \\
        \hline
        PKS-J & 601 $\pm$ 44 & Normal & Reversed & 175.6 & -32.5 & 77 & 5.5 & >773 ka (Matuyama) \\
        PKS-K & 407 $\pm$ 18 & Normal & Reversed & 203.3 & -33.5 & 18 & 12.6 & >773 ka (Matuyama) \\
        \hline
    \end{tabular}
    
    \vspace{0.5em}
    \raggedright
    \footnotesize
    Both sites show reversed polarity with good-to-excellent statistical parameters\\
    Well-defined directions argue against remagnetization\\
    Likely represent older Matuyama flows incorrectly dated\\
    Recommend new \textsuperscript{40}Ar/\textsuperscript{39}Ar dating for both sites
\end{table*}

% END OF TABLES
